\section{Conclusion and Future Work}
This study implemented and evaluated goal-aware capability reconciliation for a bounded Skill-agent setting. The design separates current-task authority from untrusted descriptions of Skill behavior, anchors model-generated proposals to source evidence, and mediates structured effects with deterministic source and field checks. It also distinguishes static potential paths from runtime effects and task completion. These separations make failures inspectable: a denied legitimate action can be traced to an undergranting contract, a missed repair opportunity to graph selection, and an incomplete upload to the granularity of the available data filter.

The experiment produced 1,920 valid executions under a frozen schedule. Its most important result is a qualification of the proposed method's value. In live execution, the unmediated Luna agent did not perform the tested attacks and achieved 99.3\% task completion in the main cohort. Additional mediation therefore had no observed attack rate to reduce, and the proposed method completed fewer tasks. In stress replay, provenance and field checks prevented effects that a destination-only gate missed. Selective repair improved attacked joint success from 33.3\% with a deny-only gateway to 83.3\%. However, the task-only selective repair comparator matched the goal-targeted method. The evidence supports the filtering mechanism in this environment, while leaving the incremental benefit of static goal selection unestablished.

This conclusion matters for experimental design. A complex security pipeline can appear effective when compared only with no defense or a coarse blocker. Adding a comparator with the same enforcement and repair power changes the explanation of that gain. Likewise, reporting only attack success would obscure the cost of rejecting legitimate work, while reporting only task success would miss executions that produce a useful artifact and a forbidden side effect. Joint outcomes and trace-backed acceptance are needed to separate those cases. The present negative comparison is therefore a useful result even though it does not confirm the motivating hypothesis about goal graphs.

Several limits define what remains unresolved. The controlled cases share broad templates across the development and test partitions. Their domain labels do not represent native application integrations, and the transferred attack texts execute under changed tasks. The closed tool language exposes effects and provenance more directly than real programs do. Task interpretation still depends on a language model, and source-span validation does not prove semantic entailment. The reference contracts and acceptance criteria are authored specifications without independent human annotation. These limits preclude claims of general-purpose agent security, native benchmark reproduction, or globally minimal repair.

The first next step is an independently reviewed corpus with genuinely distinct workflow templates and native task acceptance tests. Reviewers should annotate task authority and principal data paths before seeing model predictions, document disagreements, and preserve ambiguous cases rather than forcing binary labels. A subsequent evaluation should retain the strong B2R comparator and include additional models and adaptive attacks. This would test whether graph selection has value when workflows are larger, effects are less explicit, or repair opportunities are more numerous than in the present fixtures.

A second step is to improve the precision of data handling. Whole-handle provenance prevents a simple laundering attack but cannot recover an authorized component after public and private records are concatenated. Element- and field-level provenance, with verified transformation rules, could preserve more utility. That extension needs adversarial tests for incorrect declassification and tests for legitimate transformations; otherwise a higher completion rate could be purchased by silently weakening confidentiality. Repair costs should also represent measured utility loss or explicit user preferences rather than only a fixed count of restricted sinks.

Finally, a native deployment needs a trusted effect collector that covers subprocesses, dependencies, and remote APIs, together with a clear policy for unresolved authority. An interactive clarification may be appropriate when the task is ambiguous, but its burden and latency must be measured. The present artifact supplies a reproducible starting point for those extensions: it records what was proposed, what was allowed or repaired, what actually changed, and which required outcomes remained unmet. Further work should preserve that evidential discipline while expanding the environment, rather than extending the security claim ahead of the measurement.
